\subsection{Problem setup}
We consider the one-dimensional temperature diffusion equation
\begin{equation}
  \frac{\partial T}{\partial t} = k_z \frac{\partial^2 T}{\partial z^2},
  \qquad z \in [0, 2],\ t \in [0, 5],
\end{equation}
with true diffusivity $k_z = 0.05$, following \citet{han2026seawater}.
% NOTE: han2026seawater in references.bib is still an empty stub -- I
% cannot fabricate its author list, volume, or page numbers. This is the
% one remaining reference that has to come from you/your advisor rather
% than from a web search, since it's your own group's source paper.
We evaluate three boundary
conditions, each with a closed-form analytical solution used as ground
truth:
\begin{itemize}
  \item \textbf{Dirichlet}: $T = e^{-k_z(\pi/2)^2 t}\sin(\pi z/2) +
    e^{-k_z \pi^2 t}\sin(\pi z)$.
  \item \textbf{Neumann}: the corresponding cosine-series solution.
  \item \textbf{Robin}: $\partial T/\partial z + T = 0$ at $z=0,2$; the
    general-term solution follows \citet{han2026seawater}, Eq.~18.
\end{itemize}
For each boundary condition we solve both a \emph{forward} problem
($k_z$ known, solve for $T$) and an \emph{inverse} problem ($k_z$ unknown,
estimated jointly with $T$ from sparse observations), giving six scenarios
in total.

\subsection{Classical PINN}
The classical baseline is a fully connected network mapping $(z, t) \mapsto
T$: five hidden layers of 32 units with $\tanh$ activations. The loss
combines the PDE residual, boundary-condition residual, and (for inverse
problems) a data-fit term against sparse noisy observations; $k_z$ is a
trainable scalar in the inverse case (parameterized as $\exp(\log k_z)$ for
positivity). Training samples 2540 PDE-residual (collocation) points, 80
boundary points (40 per side), and 160 initial-condition points, each
redrawn uniformly at random per run; inverse scenarios additionally use 12
noiseless observation points on a fixed $4\times 3$ $(z,t)$ grid. The loss
term weights are $w_{bc}=w_{ic}=w_{data}=1$ throughout, and $w_{pde}=1$ for
Dirichlet and Neumann scenarios; Robin scenarios use $w_{pde}=10$, since
the gradient-dependent Robin boundary condition is harder to satisfy at
$w_{pde}=1$ and needs the PDE residual up-weighted to converge (an
asymmetry inherited from the original classical replication rather than
introduced for this study). Training is two-phase: Adam ($\mathrm{lr}=
10^{-3}$) for a fixed number of steps (5000 for Dirichlet/Neumann, 20000
for Robin, matching the higher $w_{pde}$), followed by L-BFGS (strong-Wolfe
line search, history size 100, $\mathrm{tolerance\_grad}=10^{-9}$,
$\mathrm{tolerance\_change}=10^{-11}$, up to 10000 or 20000 iterations
matching the Adam budget) until L-BFGS's own convergence tolerances are
met or the iteration budget is exhausted, whichever comes first -- there
is no early stopping based on held-out validation error.

\subsection{QCPINN}
The quantum branch is inspired by, but not identical to, the qubit-based
(discrete-variable) circuit family in \citet{farea2025qcpinn}. Inputs
$(z,t)$ are linearly rescaled to $[-1,1]$ and encoded via
\texttt{AngleEmbedding} (RX rotations) onto $n$ qubits, matching their
angle-embedding scheme. Each of $L$ variational layers applies RZ-RX
rotations on every qubit, a \emph{ring} of CNOT entangling gates (qubit $i$
to qubit $(i+1) \bmod n$, wrapping from the last qubit back to the first),
then a mirrored RX-RZ rotation layer; the circuit output is the vector of
Pauli-$Z$ expectation values $\langle Z_i\rangle \in [-1,1]$, matching
their readout. \citet{farea2025qcpinn} test four named topologies, none of
which is exactly this combination: their \emph{Cascade} topology uses ring
connectivity but continuously-tunable CRX gates rather than fixed CNOT,
while their \emph{Alternate} and \emph{Layered} topologies use fixed CNOT
but nearest-neighbor (not ring) connectivity. We therefore refer to our
circuit as \emph{ring-CNOT}: a fifth ring-plus-fixed-entangler combination
that is closest in spirit to their Cascade topology but was not itself
among the four they swept, rather than an instance of any one of their
named topologies. This output is post-processed by a small
classical head to produce $T$ (and, for inverse problems, $k_z$). Circuit
parameters are simulated exactly (no shot noise, no hardware noise) via
PennyLane's \texttt{default.qubit} device with the \texttt{backprop}
differentiation method -- \citet{farea2025qcpinn} state they use the same
shots-None/backprop configuration for their qubit-based circuits but do not
print the specific simulator device string, so this match is inferred, not
confirmed verbatim. The baseline configuration uses $n=4$ qubits and $L=2$
layers; note this differs from the $n=5$, $L=1$ configuration
\citet{farea2025qcpinn} found best after their own sweep -- Section
\ref{sec:experiments} sweeps qubit count and depth independently around our
own baseline, not theirs. The quantum branch is trained differently from
the classical one, not only architecturally but in optimizer choice: it
never switches to L-BFGS -- each L-BFGS step would require multiple
closure evaluations through the quantum circuit, which is too expensive
per step -- and instead runs Adam throughout, for a step budget equal to
the classical branch's combined Adam+L-BFGS budget (i.e.\ 15000 steps for
Dirichlet/Neumann, 40000 for Robin), with cosine learning-rate decay from
$10^{-3}$ to a floor $\eta_{\min}=10^{-5}$. This means the two branches
differ in optimizer (first-order throughout vs.\ first-order then
second-order) as well as in network type, a confound Section
\ref{sec:experiments}'s phase-1 ablation partially addresses by varying
$\eta_{\min}$ but does not eliminate -- both branches were given a matched
total step/iteration budget, not a matched optimizer.

\subsection{Training and evaluation}
All experiments use PyTorch with fixed seeds. Seeding differs by
experiment: the main classical-vs-quantum comparison
(Section~\ref{sec:experiments}, Table~\ref{tab:main-results}) uses a
single seed (0) per scenario, so the reported numbers there are point
estimates with no variance across runs; the phase-1 and phase-2 ablations
each use 3 seeds (0, 1, 2) per (scenario, configuration) pair, and it is
only for those that we report a mean and standard deviation. We return to
what this asymmetry means for how the main comparison should be read in
Section~\ref{sec:discussion}'s limitations. We report:
\begin{itemize}
  \item \textbf{MAE} and \textbf{relative $L_2$ error} against the
    analytical solution on a held-out grid;
  \item \textbf{PDE residual} at collocation points, as an
    optimization-quality check independent of the analytical solution;
  \item \textbf{$k_z$ estimation error} for inverse scenarios;
  \item \textbf{Wall-clock training time}, measured on an Apple M2 (8
    cores, 16GB RAM), CPU only for both branches -- neither PyTorch nor
    PennyLane was given GPU access, so the reported ratios are a CPU-only
    comparison. This likely does not favor the classical branch
    disproportionately: a GPU-accelerated classical network would train
    faster still, which would widen rather than narrow the gap, whereas
    the quantum branch's cost is dominated by classical simulation of the
    circuit itself and would not benefit comparably from the same
    hardware upgrade;
  \item \textbf{Circuit-only gradient norm} (phase-2 ablation only): the
    norm of the loss gradient with respect to the VQC parameters alone,
    logged at the start and end of training, isolating circuit
    trainability from the classical pre/post-processing layers.
\end{itemize}
